% Created 2025-01-28 mar 11:41
% Intended LaTeX compiler: pdflatex
\documentclass[a4paper]{article}
\usepackage[utf8]{inputenc}
\usepackage[T1]{fontenc}
\usepackage{graphicx}
\usepackage{longtable}
\usepackage{wrapfig}
\usepackage{rotating}
\usepackage[normalem]{ulem}
\usepackage{amsmath}
\usepackage{amssymb}
\usepackage{capt-of}
\usepackage{hyperref}
\usepackage[spanish]{babel}
\usepackage[usenames,dvipsnames]{color} % Required for custom colors
\renewcommand{\ttdefault}{pcr} % MONOESPACIO CON NEGRIT
\usepackage{lastpage}
\usepackage{listings}
\usepackage{listingsutf8}
\renewcommand{\lstlistingname}{Listado}
\lstset{frame=single,inputencoding=utf8,basicstyle=\scriptsize\ttfamily,showstringspaces=false,numbers=none}
\definecolor{MyDarkGreen}{rgb}{0.0,0.4,0.0} % This is the color used for comments
\lstset{ breaklines=true, postbreak=\mbox{\textcolor{red}{$\hookrightarrow$}\space}, keywordstyle=\bfseries, keywordstyle=[1]\color{Blue}\bfseries,  keywordstyle=[2]\color{Purple}\bfseries,  keywordstyle=[3]\color{Blue}\underbar,   identifierstyle=,   commentstyle=\usefont{T1}{pcr}{m}{sl}\color{MyDarkGreen}\small,   stringstyle=\color{Purple},   showstringspaces=false,   tabsize=2,   morecomment=[l][\color{Blue}]{...} }
\lstset{literate=  {á}{{\'a}}1 {é}{{\'e}}1 {í}{{\'i}}1 {ó}{{\'o}}1 {ú}{{\'u}}1   {Á}{{\'A}}1 {É}{{\'E}}1 {Í}{{\'I}}1 {Ó}{{\'O}}1 {Ú}{{\'U}}1   {à}{{\`a}}1 {è}{{\`e}}1 {ì}{{\`i}}1 {ò}{{\`o}}1 {ù}{{\`u}}1   {À}{{\`A}}1 {È}{{\'E}}1 {Ì}{{\`I}}1 {Ò}{{\`O}}1 {Ù}{{\`U}}1   {ä}{{\"a}}1 {ë}{{\"e}}1 {ï}{{\"i}}1 {ö}{{\"o}}1 {ü}{{\"u}}1   {Ä}{{\"A}}1 {Ë}{{\"E}}1 {Ï}{{\"I}}1 {Ö}{{\"O}}1 {Ü}{{\"U}}1   {â}{{\^a}}1 {ê}{{\^e}}1 {î}{{\^i}}1 {ô}{{\^o}}1 {û}{{\^u}}1   {Â}{{\^A}}1 {Ê}{{\^E}}1 {Î}{{\^I}}1 {Ô}{{\^O}}1 {Û}{{\^U}}1   {œ}{{\oe}}1 {Œ}{{\OE}}1 {æ}{{\ae}}1 {Æ}{{\AE}}1 {ß}{{\ss}}1   {ű}{{\H{u}}}1 {Ű}{{\H{U}}}1 {ő}{{\H{o}}}1 {Ő}{{\H{O}}}1   {ç}{{\c c}}1 {Ç}{{\c C}}1 {ø}{{\o}}1 {å}{{\r a}}1 {Å}{{\r A}}1   {€}{{\euro}}1 {£}{{\pounds}}1 {«}{{\guillemotleft}}1   {»}{{\guillemotright}}1 {ñ}{{\~n}}1 {Ñ}{{\~N}}1 {¿}{{?`}}1 }
\usepackage{caption}
\usepackage{attachfile2}
\usepackage[margin=1.5cm,includeheadfoot,includehead,includefoot]{geometry}
\hypersetup{colorlinks,linkcolor=black}
\usepackage{fancyhdr}
\pagestyle{fancyplain}
\chead{}
\lhead{}
\rhead{}
\cfoot{}
\lfoot{\begin{footnotesize}alvaro.gonzalezsotillo@educa.madrid.org\end{footnotesize}}
\rfoot{\begin{footnotesize}\thepage / \pageref{LastPage}\end{footnotesize}}
\usepackage[skins]{tcolorbox}
\usepackage{multicol}
\usepackage{changepage} %ajdustwidth
\usepackage{fancybox}
\usepackage{attachfile2}
\lhead{Extraordinaria 2021 (es el \\lhead)}
\rhead{Administración y Gestión de Bases de Datos (es el \\rhead)}
\lhead{IAW}
\rhead{Práctica PHP (II)}
\usepackage{svg}
\usepackage{letltxmacro}
\LetLtxMacro{\originalincludegraphics}{\includegraphics}
\renewcommand{\includegraphics}[2][]{\IfFileExists{#2.pdf}{\originalincludegraphics[#1]{#2.pdf}}{\originalincludegraphics[#1]{#2}}}
\LetLtxMacro{\originalincludesvg}{\includesvg}
\renewcommand{\includesvg}[2][]{\IfFileExists{#2.pdf}{\originalincludegraphics[#1]{#2.pdf}}{\originalincludegraphics[#1]{#2.svg.pdf}}}
\usepackage{comment}
\excludecomment{NOTES}
\date{}
\title{Tercera práctica PHP}
\hypersetup{
 pdfauthor={Álvaro González Sotillo},
 pdftitle={Tercera práctica PHP},
 pdfkeywords={},
 pdfsubject={},
 pdfcreator={Emacs 29.4 (Org mode 9.6.15)}, 
 pdflang={Spanish}}
\begin{document}

\maketitle
\setcounter{tocdepth}{1}
\tableofcontents

\captionsetup{font=scriptsize}

\textattachfile{/home/alvaro/apuntes-clase/implantacion-aplicaciones-web-asir2/apuntes/UT05/UT05-practica-php-ajax.org}{} \vspace{-1em}


\setlength{\parindent}{0em}
\setlength{\parskip}{1em}

\newtcolorbox{Aviso}[1][Aviso]{
  enhanced,
  colback=gray!5!white,
  colframe=gray!75!black,fonttitle=\bfseries,
  colbacktitle=gray!85!black,
  attach boxed title to top left={yshift=-2mm,xshift=2mm},
  title=#1
}

\newtcolorbox{cuadrito}[1][Ignorado]{
  %drop shadow southeast,
  enhanced jigsaw,
  colback=white,
}


\newcommand{\StudentData}{
  \begin{cuadrito}[1\textwidth]
    \vspace{0.3cm}
    \large{
      \textbf{Apellidos:} \hrulefill \\
      \textbf{Nombre:} \hrulefill \\
      \textbf{Fecha:} \hrulefill \hspace{1cm} \textbf{Usuario:} \hrulefill \\
    }
    \vspace{-0.2cm}
  \end{cuadrito}
}

\newcommand{\StudentDataDniFirma}{
  \begin{cuadrito}[1\textwidth]
    \vspace{0.3cm}
    \large{
      \textbf{Apellidos:} \hrulefill \\
      \textbf{Nombre:} \hrulefill  \\
      \textbf{Dni:} \hrulefill \hspace{8cm} \\
      \\
      \textbf{Firma:} \hrulefill \hspace{8cm} \\
    }
    \vspace{-0.2cm}
  \end{cuadrito}
}

\attachfile{../UT04//tiles/tiles.zip}

\section*{Objetivos de la práctica:}
\label{sec:org0000000}
En esta práctica se espera que el alumno:
\begin{itemize}
\item Se familiarice con el manejo de sesiones
\item Se familiarice con el uso de MySQL
\item Utilice técnicas AJAX para mejorar la usabilidad de sus aplicaciones web
\end{itemize}


Se tomará como base el código desarrollado para la práctica anterior

\section{\emph{CRUD} de mazmorras (3 puntos)}
\label{sec:org0000003}
Modifica \texttt{listado-mazmorras.php}
\begin{itemize}
\item El botón de crear una nueva mazmorra solo podrá pulsarse si el campo de texto no está vacío, y el nombre de mazmorra aún no existe.
\item hay que tener en cuenta que otros usuarios pueden estar creando mazmorras
\begin{itemize}
\item Debe ser una petición AJAX, no puede descargarse una lista de nombres prohibidos al comienzo de la página
\end{itemize}
\item El borrado de mazmorra no recargará la página
\begin{itemize}
\item La fila de la tabla aparecerá tachada, y los botones deshabilitados
\item Si se recarga la página, entonces sí que ya no aparecerá la fila
\end{itemize}
\end{itemize}

\section{Duplicar mazmorra (1 punto)}
\label{sec:org0000006}
Modifica \texttt{duplicar-mazmorra.php}. Está página
\begin{itemize}
\item El botón de duplicar mazmorra solo se podrá pulsar si el nuevo nombre no está vacío, y el nombre e mazmorra aún no existe.
\end{itemize}

\section{Editar mazmorra  (5 puntos)}
\label{sec:org0000009}
Modifica \texttt{editar-mazmorra.php}. Esta página:
\begin{itemize}
\item Inicialmente, estará seleccionada la imagen \texttt{E.png} en la paleta
\item Si se hace click sobre otro item de la paleta, se seleccionará
\item Bastará con hacer click en la mazmorra para sustituir esa tesela por la seleccionada en la paleta.
\end{itemize}

\section*{Recomendaciones}
\label{sec:org000000c}
\begin{itemize}
\item Crea una página o varias páginas (por ejemplo \texttt{api-mazmorras.php}) donde se realizarán las peticiones AJAX
\item Decide si utilizarás JSON o POST para estas páginas de API
\end{itemize}


\section*{Normas de entrega}
\label{sec:org000000f}
\begin{itemize}
\item El código estará en el directorio \texttt{\textasciitilde{}/public-html/practica03/} del \emph{workspace} de Coder del alumno
\begin{itemize}
\item El profesor copiará las páginas a su propio LAMP para corregir la práctica
\item El profesor se encargará de cambiar los datos de acceso a MySQL.
\item Las páginas deberán crear las tablas que necesiten, si no no se podrá corregir la práctica.
\end{itemize}
\item Los nombres de fichero PHP y los nombres de funciones serán los especificados en los enunciados
\item Todos los formularios serán \textbf{GET}
\begin{itemize}
\item Si no se pueden ver los parámetros en la URL, el ejercicio no tendrá más del 50\% de la nota
\end{itemize}
\item Se recuerda que los alumnos deben guardar una copia de seguridad de todos sus trabajos, en el caso de que Coder no funcione
\item \textbf{No deben aparcer errores ni avisos de PHP}. Hay que tener en cuenta que se ejecutarán con opciones de \emph{debug}, como \texttt{display\_errors} a \texttt{true}
\begin{itemize}
\item Puede ser útil: \texttt{isset()}, \texttt{is\_numeric()}, \texttt{floatval()}
\item Si hay errores o avisos, el ejercicio no tendrá más del 50\% de la nota
\end{itemize}
\item No se tendrán en cuenta:
\begin{itemize}
\item Errores debidos al reenvío de formularios por duplicado (por recarga de página)
\item Errores debidos a parámetros cambiados a mano para teselas/mazmorras que no existen
\item Errores debidos a excesiva latencia de la respuesta del servidor
\end{itemize}
\end{itemize}

\section*{Resultados de aprendizaje}
\label{sec:org0000012}
Esta práctica contribuye a la calificación de los siguientes RA:
\begin{itemize}
\item RA 5: parte de "prácticas"
\item RA 6: parte de "prácticas"
\end{itemize}
\end{document}
